\section{What the Router Caches}
\label{sec:ch18-what-the-router-caches}

\index{caching!at the router|(}%
Sections~\ref{sec:ch18-writing-it-down} and~\ref{sec:ch18-when-it-fails-instead}
leave every service able to say what a request cost it and why it failed. The usual reply to a slow page, before anyone has
found the service, is to put a cache in front of it, and the router is the
obvious place: every request passes through it. So it is worth knowing
exactly what this router caches and what that saves, in the counts this book
uses.

\subsection{The Plan, Not the Answer}

\index{query plan!cached, and what that saves}%
The router caches its own work. Normalizing an operation and planning its fetches are done once per distinct
operation and reused.

\code{enable\_cache\_response\_headers}, under \code{engine.debug}, makes the
router report each of those caches as a hit or a miss in a response header of
its own~\autocite{cosmorouterconfig}. Each router setting in this section is appended to \code{config.yaml} for one
run, the router is restarted on it, and the book's own \code{config.yaml} is
put back afterwards; none of them ships. With this one on, send the page from
section~\ref{sec:ch18-the-page-that-got-slower} twice, on the graph as this
chapter leaves it, and count statements the way
section~\ref{sec:ch18-writing-it-down} does:

\begin{center}
\begin{tabular}{lcrrrr}
\toprule
Request & Plan cache & Sessions & Speakers & Ratings & Search \\
\midrule
First & \code{MISS} & 1 & 1 & 1 & 0 \\
Second & \code{HIT} & 1 & 1 & 1 & 0 \\
\bottomrule
\end{tabular}
\end{center}

\index{normalization!a literal lifted into a variable shares the plan}%
The second request skipped planning and every service did exactly the work it
did the first time. That is not a weakness of the cache; it caches a plan,
and a plan is instructions for asking the services, not their answers. It is
also shared more widely than the query text suggests. Ask for
\code{sessions(first: 2)} instead of ten and the normalization cache misses
while the plan cache hits. Chapter~\ref{ch:enter-the-router} showed the router
lifting a literal argument into a variable before it forwards an operation,
and here the two page sizes share one plan. None of it changes how many statements the Speakers
service spends, so none of it would have made
section~\ref{sec:ch18-the-page-that-got-slower}'s page faster.

\subsection{A Header for Somebody Else's Cache}

\index{Cache-Control@\code{Cache-Control}!computed by the router|(}%
What the router can do about answers is tell a cache in front of it how long
one may be kept. \code{cache\_control\_policy} is off by default. Turned on, it
writes a \code{Cache-Control} header on each response from the values the
subgraphs involved in that request carry, taking \code{no-store} or
\code{no-cache} from any of them over everything else and otherwise the
smallest \code{max-age}; a subgraph's value is either the \code{Cache-Control}
header it sends back or one set for it in the router's config; and a response
with errors in it gets \code{no-store} whatever the subgraphs
said~\autocite{cosmocachecontrol}. None of this book's services sends a \code{Cache-Control} of its own, so here
every value comes from the router's config.

\index{Cache-Control@\code{Cache-Control}!the most restrictive subgraph decides}%
Restart the router with the policy enabled, a global \code{value} of
\code{max-age=180, public}, and a \code{subgraphs} entry giving \code{ratings}
\code{max-age=60, public}. The header then depends on which services a request
happened to reach. The session
titles and speakers come back cacheable for 180 seconds. Add
the two rating fields, which only the Ratings service answers, and the whole
page is cacheable for 60. The value belongs to the response, so the least cacheable
service a request touches decides for every field in it, the same shape as the
blast radius chapter~\ref{ch:null-blast-radius} measured for nulls.

\index{Cache-Control@\code{Cache-Control}!a subgraph with no value has no say}%
A service with no value set has no say at all. Restart the router
with the global value taken away and only \code{speakers} given a value,
\code{max-age=300, public}. A page of titles, \code{\{ sessions(first: 10) \{ nodes \{ title \} \} \}},
which is all Sessions data, gets no \code{Cache-Control} header. The same page
with speaker names on it gets \code{public} and 300 seconds, although every
title on it came from a service that expressed no view on how long it may be
kept. The router reads a missing value as no restriction at all.

\index{Cache-Control@\code{Cache-Control}!public on an authorized response}%
And the policy does not know who asked. Restart the router with the global value back at
\code{max-age=180, public}, send the session page asking for each speaker's email, the column
chapter~\ref{ch:entities-authorization} guarded:

\begin{minted}{graphql}
{ sessions(first: 10) { nodes { title speaker { name email } } } }
\end{minted}

 Without a token the email is refused,
the response has errors, and it is marked \code{no-store}. With the bearer
token, the guarded column comes back, and the response carrying it is marked
\code{public} for 180 seconds with nothing in \code{Vary} beyond
\code{Accept-Encoding}. RFC 9111 is specific about what that permits: a shared
cache must not reuse a response to a request that carried an
\code{Authorization} header unless the response carries a directive allowing
it, and \code{public} is one of the three that do~\autocite{rfc9111}.
\index{Cache-Control@\code{Cache-Control}!computed by the router|)}%

\subsection{The Cache This Router Does Not Have}

\index{response cache!not in the router this book pins}%
WunderGraph's documentation describes a response cache, one that stores the
subgraph answers themselves, entity fetches included, keyed by entity and
selection set and governed by each subgraph's \code{Cache-Control}. The page
marks it experimental, \enquote{in alpha} and not to be used in production, and disabled by default~\autocite{cosmoresponsecache}. The router this book runs does not
have it: add a \code{response\_cache} block to \code{config.yaml} and the router
refuses to start, naming it as a property its configuration schema does not
allow, the same refusal chapter~\ref{ch:enter-the-router} got for a key it made
up. Whatever release the page describes, it is not the one
appendix~\ref{app:toolchain} pins.

\index{caching!my judgment on the router}%
So this book ships no cache policy, and that is my judgment rather than
something measured. The router this book runs has two levers: the plans it
caches, which it caches already and which save no service anything, and the
length of time it tells a shared cache to keep a whole response, which is decided per response out of
values set per service. The reason this graph must not be cached is one column
in one service, behind a token. A per-response setting built from per-service
values cannot express that, and the one it will express quietly is
\code{public}. When a page is slow, the id in
section~\ref{sec:ch18-writing-it-down} finds the service that made it slow, and
the fix belongs in that service.
\index{caching!at the router|)}%
